{
\titleformat
	{\subsection} % command
	[block] % shape
	{\bfseries\large} % format
	{} % label
	{0.2ex} % sep
	{
	    \begin{minipage}[t]{\linewidth}
	    	\raggedright \nohyphenation{#1}
	    \end{minipage}
	} % before-code
	[% after code
		\vspace{-1.5ex}
	]

\GRPchapter{Races}
Hi hi hiiiiiii

Here is the list of RAAACES. they aren't RACIST they are just RACES. YEEEEEAH.

Species provide either of

\begin{itemize}
	\item \GRPdicebonus{2} to two different talents
	\item \GRPdicebonus{2} to one talent and \GRPdicebonus{1} to two different talents.
	\item \GRPdicebonus{3} to one talent
\end{itemize}

\GRPpagebreak
\section{Wyrmkin}

The wyrmkin are the most populous race in the galaxy (though by themselves they do not comprise a majority of sentients). They have many dragon-like qualities, including their scales, their elemental breath, and their hatching from warm incubated eggs. It is a common belief among the wyrmkin that they carry the blood of the ancient dragons within them because of these similarities. 

The wyrmkin have many names for themselves, the most common being ``shell-born'' -- referring to their hatching from eggs which is uncommon among sentient races. Other common aliases include ``of the holy breath'' and ``chosen of the dragons,'' but these days these phrases are typically only heard among groups entrenched in an attitude of wyrmkin superiority over other races.

\subsection*{Appearance}


\subsection*{Prehensile Limbs}

Wyrmkin are of humanoid form with a non-prehensile tail.

\GRPstat{Grab Limbs} is 2.

\subsection*{Languages}

\begin{itemize}
	\item Standard
\end{itemize}

\subsection*{Species Bonuses}

\begin{itemize}
	\item \GRPdicebonus{2}\GRPtalent{Intellect}
	\item \GRPdicebonus{2}\GRPtalent{Charisma}
\end{itemize}

\subsection*{Species Abilities}

\begin{itemize}
	\item \textbf{Wyrm Breath}-- Choose one of the \emph{transformative elements} (Fire or lightning). As an action you may produce this element from your mouth in a controlled fashion. If using this to craft or create something, relevant talent checks must still be made. During \emph{phased} mode, you may use this to fill a room temporarily with the element of your choice as a Room Range attack, using Fortitude for your attack roll.
	\item \textbf{Draconic Upbringing}-- Start with 1 rank in the skill \GRPability{Familiar with Dragon Culture} at the creation of your character before investing any experience. If this does not apply to the setting you are using this character in, then you may replace this with a different \GRPability{Familiar} skill at the discretion of the GM.
\end{itemize}

\GRPpagebreak
\section{Kobolds}

Kobolds look superficially similar to the wyrmkin, but have much smaller proportions and have a more raptor-like stature. Unlike the wyrmkin, kobolds have live births rather than hatching from incubated eggs. More insular wyrmkin tend to look down on kobolds as a parody of their race. Historically, kobold communities have faced considerable prejudice from wyrmkin society.

\subsection*{Appearance}

Kobolds are covered with scales like many of the reptilian races, but may also have small amounts of featherlike softer scales in some areas of their bodies. Kobold scale colors tend to be either shades of orange, brown, blue, or green, and may have spotted and striped patterns.

\subsection*{Prehensile Limbs}

Kobolds are bipedal with two prehensile arms and a prehensile tail.

\GRPstat{Grab Limbs} is 3.

\subsection*{Languages}

\begin{itemize}
	\item Standard
	\item Slewbask
\end{itemize}

\subsection*{Species Bonuses}

\begin{itemize}
	\item \GRPdicebonus{1}\GRPtalent{Memory}
	\item \GRPdicebonus{2}\GRPtalent{Insight}
	\item \GRPdicebonus{1}\GRPtalent{Dexterity}
\end{itemize}

\subsection*{Species Abilities}

\begin{itemize}
	\item \textbf{Pack Instinct}-- During offensive phased mode, making any sort of helping roll does not consume an action. If helping make a cast or any other consumable ability, a kobold still consumes a cast use or the consumable resource.
\end{itemize}

\GRPpagebreak
\section{Snakefolk}

\subsection*{Appearance}

\subsection*{Prehensile Limbs}

\GRPstat{Grab Limbs} is 2

\subsection*{Languages}

\subsection*{Species Bonuses}

\begin{itemize}
	\item \GRPbonusdice{2}\GRPtalent{Athletics}
	\item \GRPbonusdice{1}\GRPtalent{Will}
	\item \GRPbonusdice{1}\GRPtalent{Charisma}
\end{itemize}

\subsection*{Species Abilities}


\GRPpagebreak
\section{Carcinids}

Carcinids are large, crab-like sentient beings that live in all corners of the galaxy. Many independent communities of carcinids survived the Spiral War in many different.

Most carcinids do not have anything that could be denoted a biological sex. Most carcinids may reproduce by laying eggs which may contain just their own genetic material, or include that of a co-parent.

\subsection*{Appearance}

Carcinids have a hard outer carapace around a scaffolding they use as an exoskeleton. They are of smaller stature than most sentients. There are many possible colorations and patterns for a carcinid's carapace.

\subsection*{Prehensile Limbs}

Carcinids have 10 limbs they can walk on, 5 of which are free for manipulation of objects at any given time.

\GRPstat{Grab Limbs} is 5

\subsection*{Languages}
Carcinids are unable to learn to speak any language besides chitter without the use of assistive devices or spells due to the shape of their mouthparts.

\begin{itemize}
	\item Standard (limited)
	\item Chitter
\end{itemize}

\subsection*{Species Bonuses}

\begin{itemize}
	\item \GRPdicebonus{2}\GRPtalent{Dexterity}
	\item \GRPdicebonus{2}\GRPtalent{Insight}
\end{itemize}

\subsection*{Species Abilities}

\begin{itemize}
	\item 
\end{itemize}

\GRPpagebreak
\section{Automa}

Automa are living machines from the Golden Age. All of them were shutdown by the same weapon that destroyed the fey corridor network. Old broken down automa are often found in Golden Age ruins, and occasionally may be in a state where they may be restored to operation. Once restored to operation, they are able to heal from injuries and 

Automa tend to have very degraded and vague memories of their time before the Spiral War, but this is still considered valuable information by many artifact hunters.

\subsection*{Appearance}

\subsection*{Prehensile Limbs}

Automa come in various shapes and have a minimum of 2 prehensile limbs. They may have up to 2 additional prehensile limbs determined by the \GRPability{Configurable Body} species ability.

\subsection*{Languages}
\begin{itemize}
	\item Standard
\end{itemize}

\subsection*{Species Bonuses}
\begin{itemize}
	\item \GRPdicebonus{2}\GRPtalent{Intellect} - during the \GRPability{Recuperate} action, you can forego this species bonus to have one additional prehensile limb
	\item \GRPdicebonus{2}\GRPtalent{Fortitude} -  during the \GRPability{Recuperate} action, you can forego this species bonus to have one additional prehensile limb
\end{itemize}

\subsection*{Species Abilities}
\begin{itemize}
	\item \textbf{Essentia Essentials}-- Automa do not need to eat or breathe. They appear to draw on the same power that Fey Corridors draw on. When taking the \GRPability{Recuperate} downtime action, Automa do not consume rations. Additionally they are unaffected by hostile atmospheres or a lack of atmosphere. 
	\item \textbf{Configurable body}-- Whenever you take the \GRPability{Recuperate} downtime action, you may add up to 2 additional \GRPstat{Grab Limbs} by forgoing your Species bonuses.

	Each Species bonus you forgo gives you 1 additional Grab Limb.
	\item \textbf{}
\end{itemize}

\GRPpagebreak
\section{Species template}

\subsection*{Appearance}

\subsection*{Prehensile Limbs}

\subsection*{Languages}

\subsection*{Species Bonuses}

\subsection*{Species Abilities}

\closeopenenv{multicols}
}